\documentclass[11pt]{article}
\usepackage[a4paper,margin=1in]{geometry}
\usepackage{amsmath,amssymb,mathtools}
\usepackage{physics}
\usepackage{siunitx}
\usepackage{graphicx}
\usepackage{bm}
\setlength{\parskip}{0.6em}
\setlength{\parindent}{0pt}

\begin{document}

\begin{center}
{\Large HW2 Problem 1a  Derivation of the beam model with a point load and self weight}\\
\end{center}

\textbf{Kinematics and constitutive law.}
Let $\phi(x)$ be the small transverse deflection of a slender uniform beam of length $L$ resting in the $x$ direction.
Euler Bernoulli assumptions are used
plane cross sections remain plane and normal to the neutral axis
shear strains are negligible
and rotations are small.
The curvature of the neutral axis is
\[
\kappa(x) \approx \dv[2]{\phi}{x}.
\]
For an elastic beam with Young modulus $E$ and second moment of area $I$ the bending moment is related to curvature by
\[
M(x) = - E I \, \kappa(x) = - E I \, \phi''(x).
\]
The sign choice is such that positive downward deflection produces negative bending moment.

\textbf{Force balance in the beam.}
Let $q(x)$ be the distributed transverse load per unit length taken positive downward.
Classical beam theory gives the static relations
\[
V(x) = \dv{M}{x}, 
\qquad 
q(x) = \dv{V}{x} = \dv[2]{M}{x}.
\]
Combining with the constitutive law yields the governing equation
\[
E I \, \phi^{(4)}(x) = q(x).
\]

\textbf{Loads in the present setting.}
Two sources act on the beam
\begin{itemize}
\item a concentrated force of magnitude $P$ applied at $x = L/3$ directed downward
\item uniform self weight of the beam with total mass $m$ and gravity $g$.  
Since the mass is uniform the weight per unit length is $\omega = \dfrac{m g}{L}$.
\end{itemize}
In distribution form the total load is
\[
q(x) = P\,\delta\!\left(x-\frac{L}{3}\right) + \omega
= P\,\delta\!\left(x-\frac{L}{3}\right) + \frac{m g}{L}.
\]
Therefore the field equation reads
\[
E I \, \phi^{(4)}(x) 
= P\,\delta\!\left(x-\frac{L}{3}\right) + \frac{m g}{L}.
\]

\textbf{Boundary conditions for clamped ends.}
Both ends are fixed horizontally and vertically.  
For a clamped end the deflection and the rotation vanish, which gives
\[
\phi(0)=0, 
\qquad 
\phi'(0)=0,
\qquad
\phi(L)=0,
\qquad
\phi'(L)=0.
\]

\textbf{Consistency checks.}
Integrating the field equation across a small interval that encloses $x=L/3$ gives the jump in the third derivative
\[
\phi^{(3)}\bigl((L/3)^{+}\bigr) - \phi^{(3)}\bigl((L/3)^{-}\bigr)
= \frac{P}{E I},
\]
which is the expected shear force jump caused by a concentrated load.
Away from the point $x=L/3$ the right side is constant, hence $M'' = q = m g L^{-1}$ and the moment is a cubic in $x$, again consistent with elementary beam theory.

\textbf{Final boundary value problem.}
With the choice $\phi$ for deflection the model used in the later parts of the assignment is
\[
\boxed{ \;
E I \, \phi^{(4)}(x) 
= P\,\delta\!\left(x-\frac{L}{3}\right) + \frac{m g}{L},
\quad
\phi(0)=\phi'(0)=0,
\quad
\phi(L)=\phi'(L)=0
\;}
\]
which is equivalent to the statement in the prompt after the change of notation $\phi \leftrightarrow v$.
\end{document}
